% GENERATED FILE. Edit canonical lessons, then run npm run generate:books.
% canonical-source-hash: fnv1a64:41e65334907f04cf
% canonical-lessons: FA-C237-barfi, FA-C237-vahshi, FA-C237-ahli, FA-C237-zende, FA-C237-morde

\chapter{Snowy, Wild, Tame, Alive, Dead}
\label{ch:fa-snowy-wild-tame-alive-dead}
\hlchaptermodality{237}

\begin{chapteropening}
\textbf{By the end of this chapter:} \emph{I can say snowy, wild, tame, alive and dead.}

Complete the last lesson of chapter 237: I can say snowy, wild, tame, alive and dead.
\end{chapteropening}

\section[\texorpdfstring{\fa{برفی} (barfi)}{برفی (barfi)}]{\fa{برفی} (barfi) — snowy}

\label{lesson:FA-C237-barfi}



\begin{quote}
Before the new one: say the Persian for crazy, then the Persian for loyal.
\end{quote}

\subsection*{You'll want to know: \fa{برفی}}
\textbf{\fa{برفی}} — \emph{barfi} — \textquotedblleft{}snowy\textquotedblright{}.

\begin{cousinweb}
\emph{barfi} is \emph{barf}, \textquotedblleft{}snow\textquotedblright{}, with the ending \emph{-i}. \emph{barf} is inherited, from Middle Persian \emph{wafr} (EtymDB).
\end{cousinweb}

\begin{grammarlens}[title={What you've built}]
One more word for this chapter.
\end{grammarlens}

\subsection*{Your turn}
\begin{itemize}
\raggedright
  \item \emph{Say it:} \emph{barfi}
  \item \emph{Say it:} \emph{barfi}, once more
  \item \emph{Say it:} say \emph{barfi}
  \item \emph{Recall:} say the Persian for crazy, then the Persian for loyal, then say \emph{barfi} again
\end{itemize}

\begin{tcolorbox}[breakable,colback=teal!4,colframe=teal!35!black,title={Before you move on}]
What does \fa{برفی} mean? (\textquotedblleft{}Snowy\textquotedblright{}.) Say it once more.
\end{tcolorbox}
\section[\texorpdfstring{\fa{وحشی} (vahshi)}{وحشی (vahshi)}]{\fa{وحشی} (vahshi) — wild}

\label{lesson:FA-C237-vahshi}



\begin{quote}
Before the new one: say the Persian for loyal, then the Persian for snowy.
\end{quote}

\subsection*{You'll want to know: \fa{وحشی}}
\textbf{\fa{وحشی}} — \emph{vahshi} — \textquotedblleft{}wild\textquotedblright{}.

\begin{cousinweb}
\emph{vahshi} is borrowed from Arabic \emph{waḥšī}, \textquotedblleft{}wild, of the wilderness\textquotedblright{}. Its Arabic root is \emph{w-ḥ-š} (Lane).
\end{cousinweb}

\begin{grammarlens}[title={What you've built}]
One more word for this chapter.
\end{grammarlens}

\subsection*{Your turn}
\begin{itemize}
\raggedright
  \item \emph{Say it:} \emph{vahshi}
  \item \emph{Say it:} \emph{vahshi}, once more
  \item \emph{Say it:} say \emph{vahshi}
  \item \emph{Recall:} say the Persian for loyal, then the Persian for snowy, then say \emph{vahshi} again
\end{itemize}

\begin{tcolorbox}[breakable,colback=teal!4,colframe=teal!35!black,title={Before you move on}]
What does \fa{وحشی} mean? (\textquotedblleft{}Wild\textquotedblright{}.) Say it once more.
\end{tcolorbox}
\section[\texorpdfstring{\fa{اهلی} (ahli)}{اهلی (ahli)}]{\fa{اهلی} (ahli) — tame, domestic}

\label{lesson:FA-C237-ahli}



\begin{quote}
Before the new one: say the Persian for snowy, then the Persian for wild.
\end{quote}

\subsection*{You'll want to know: \fa{اهلی}}
\textbf{\fa{اهلی}} — \emph{ahli} — \textquotedblleft{}tame, domestic\textquotedblright{}.

\begin{cousinweb}
\emph{ahli} is borrowed from Arabic \emph{ahlī}, \textquotedblleft{}of the household\textquotedblright{}. Its Arabic root is \emph{ʾ-h-l} (Lane).
\end{cousinweb}

\begin{grammarlens}[title={What you've built}]
One more word for this chapter.
\end{grammarlens}

\subsection*{Your turn}
\begin{itemize}
\raggedright
  \item \emph{Say it:} \emph{ahli}
  \item \emph{Say it:} \emph{ahli}, once more
  \item \emph{Say it:} say \emph{ahli}
  \item \emph{Recall:} say the Persian for snowy, then the Persian for wild, then say \emph{ahli} again
\end{itemize}

\begin{tcolorbox}[breakable,colback=teal!4,colframe=teal!35!black,title={Before you move on}]
What does \fa{اهلی} mean? (\textquotedblleft{}Tame, domestic\textquotedblright{}.) Say it once more.
\end{tcolorbox}
\section[\texorpdfstring{\fa{زنده} (zende)}{زنده (zende)}]{\fa{زنده} (zende) — alive}

\label{lesson:FA-C237-zende}



\begin{quote}
Before the new one: say the Persian for wild, then the Persian for tame.
\end{quote}

\subsection*{You'll want to know: \fa{زنده}}
\textbf{\fa{زنده}} — \emph{zende} — \textquotedblleft{}alive\textquotedblright{}.

\begin{cousinweb}
\emph{zende} is a native Persian word, inherited rather than borrowed: it continues Middle Persian \emph{zīndag}, \textquotedblleft{}alive\textquotedblright{} (EtymDB).
\end{cousinweb}

\begin{grammarlens}[title={What you've built}]
One more word for this chapter.
\end{grammarlens}

\subsection*{Your turn}
\begin{itemize}
\raggedright
  \item \emph{Say it:} \emph{zende}
  \item \emph{Say it:} \emph{zende}, once more
  \item \emph{Say it:} say \emph{zende}
  \item \emph{Recall:} say the Persian for wild, then the Persian for tame, then say \emph{zende} again
\end{itemize}

\begin{tcolorbox}[breakable,colback=teal!4,colframe=teal!35!black,title={Before you move on}]
What does \fa{زنده} mean? (\textquotedblleft{}Alive\textquotedblright{}.) Say it once more.
\end{tcolorbox}
\section[\texorpdfstring{\fa{مرده} (morde)}{مرده (morde)}]{\fa{مرده} (morde) — dead}

\label{lesson:FA-C237-morde}



\begin{quote}
Before the new one: say the Persian for tame, then the Persian for alive.
\end{quote}

\subsection*{You'll want to know: \fa{مرده}}
\textbf{\fa{مرده}} — \emph{morde} — \textquotedblleft{}dead\textquotedblright{}.

\begin{cousinweb}
\emph{morde} is the past participle of \emph{mordan}, \textquotedblleft{}to die\textquotedblright{}: \textquotedblleft{}died, dead\textquotedblright{}. The verb is native, usually traced to Middle Persian \emph{murdan}.
\end{cousinweb}

\begin{grammarlens}[title={What you've built}]
One more word for this chapter.
\end{grammarlens}

\subsection*{Your turn}
\begin{itemize}
\raggedright
  \item \emph{Say it:} \emph{morde}
  \item \emph{Say it:} \emph{morde}, once more
  \item \emph{Say it:} say \emph{morde}
  \item \emph{Recall:} say the Persian for tame, then the Persian for alive, then say \emph{morde} again
\end{itemize}

\begin{tcolorbox}[breakable,colback=teal!4,colframe=teal!35!black,title={Before you move on}]
What does \fa{مرده} mean? (\textquotedblleft{}Dead\textquotedblright{}.) Say it once more.
\end{tcolorbox}
